\chapter{Tổng quan}
\label{chap:overview}

\section{Lý do chọn đề tài}
Trong kỷ nguyên số hóa mạnh mẽ với nhịp sống hiện đại ngày nay, việc quản lý thời gian, lịch trình và công việc hằng ngày đã trở thành một thách thức lớn đối với mỗi cá nhân khi khối lượng thông tin và tác vụ cần xử lý ngày càng phức tạp. Các giải pháp quản lý hiện nay phần lớn vẫn hoạt động một cách độc lập và thụ động. Người dùng thường phải tự thao tác thủ công để đồng bộ lịch trình giữa các nền tảng khác nhau (như Google Workspace hay Lark Suite), trong khi hệ thống chỉ đóng vai trò lưu trữ tĩnh. Thực tế này mở ra nhu cầu sử dụng trợ lý ảo có khả năng tự động hóa các thao tác trên, tương tác với người dùng thông qua ngôn ngữ tự nhiên. Song song với đó, sự bùng nổ của các Mô hình Ngôn ngữ Lớn (Large Language Models - LLM) đã mở ra một chương mới cho ngành công nghệ phần mềm. Điều này dẫn đến sự xuất hiện của các tác nhân trí tuệ nhân tạo (AI Agents) với khả năng lập luận logic mạnh mẽ, hiểu sâu ngữ cảnh và linh hoạt sử dụng công cụ để thực thi tác vụ \cite{wang2023surveyagents}. Sự phát triển đột phá này chính là lời giải cho bài toán quản lý hiệu quả nêu trên. Xuất phát từ thực tiễn công nghệ và nhu cầu cấp thiết đó, đề tài “Tích hợp LLM để phát triển trợ lý ảo giúp quản lý lịch trình và công việc hằng ngày” được lựa chọn thực hiện, nhằm mục tiêu hiện thực hóa một tác nhân cá nhân thông minh, chủ động và tối ưu hóa trải nghiệm quản lý công việc của người dùng.

\section{Mục tiêu đề tài}

\subsection{Xây dựng trợ lý ảo cá nhân tự động hóa công việc}
Mục tiêu chính là xây dựng một trợ lý cá nhân hoạt động dưới dạng dịch vụ chạy nền liên tục (daemon), tự động hóa các tác vụ quản lý lịch trình hằng ngày qua giao diện Telegram:
\begin{itemize}
  \item Hệ thống tiếp nhận và xử lý yêu cầu qua tin nhắn Telegram tức thời.
  \item Thiết lập giao diện tương tác trực quan, đảm bảo thời gian phản hồi nhanh chóng.
  \item Tận dụng khả năng hiểu hội thoại tự nhiên của LLM để nâng cao trải nghiệm người dùng.
  \item Chủ động theo dõi tiến độ công việc và phát thông báo nhắc nhở khi đến hạn.
\end{itemize}

\subsection{Tích hợp tác nhân AI và thiết kế khung công cụ (Tool Harness)}
Đề tài nghiên cứu cách thức xây dựng vòng lặp tác nhân (agent loop) dựa trên ReAct kết hợp thiết kế khung công cụ (\texttt{Tool Harness}):
\begin{itemize}
  \item Chuẩn hóa lược đồ mô tả công cụ bằng JSON Schema giúp LLM gọi hàm chính xác.
  \item Triển khai khung công cụ (\texttt{Tool Harness}) hỗ trợ thực thi câu lệnh shell và CLI của các hệ thống văn phòng như Google Workspace (qua \texttt{gws}) hay Lark Suite (qua \texttt{lark-cli}) thông qua công cụ thực thi lệnh hệ thống \texttt{execute} \cite{googleworkspaceCli,larkCli}.
  \item Thiết lập tầng kiểm soát và phân quyền an toàn, ngăn chặn các hành vi leo thang đặc quyền hoặc truy cập trái phép.
\end{itemize}

\subsection{Nghiên cứu tối ưu ngữ cảnh và cơ chế nén ngữ cảnh tự động}
Giải quyết giới hạn cửa sổ ngữ cảnh và chi phí token của LLM thông qua các kỹ thuật xử lý ngữ cảnh:
\begin{itemize}
  \item Phân tầng hệ thống bộ nhớ thành hai lớp: bộ nhớ ngắn hạn duy trì dữ liệu tương tác tức thời và bộ nhớ dài hạn lưu trữ thông tin bền vững.
  \item Phát triển thuật toán nén ngữ cảnh (context compaction) tự động nhằm lọc bỏ thông tin nhiễu, chỉ lưu trữ các sự kiện và trạng thái quan trọng giúp tối ưu hóa ngân sách token và duy trì hiệu năng của tác nhân.
\end{itemize}

\subsection{Trích xuất quy trình làm việc có kiểm soát (HITL)}
Giúp tác nhân tự mở rộng năng lực hành vi qua quá trình sử dụng:
\begin{itemize}
  \item Nhận diện các mẫu hành vi lặp lại từ lịch sử tương tác của người dùng.
  \item Đề xuất đóng gói chuỗi hành động thành quy trình làm việc mới dưới dạng tệp \texttt{SKILL.md}.
  \item Áp dụng cơ chế kiểm soát bởi người dùng (HITL) để xác nhận trước khi nạp động quy trình vào hệ thống.
\end{itemize}

\subsection{Rèn luyện tư duy hệ thống và kỹ năng giải quyết vấn đề thực tiễn}
Thông qua quá trình xây dựng sản phẩm thực tế, đồ án giúp nâng cao các kỹ năng chuyên môn:
\begin{itemize}
  \item Sử dụng Node.js, TypeScript và thư viện LangGraph để thiết kế kiến trúc đồ thị trạng thái bền vững \cite{nodejsDocs,typescriptDocs,langgraphDocs}.
  \item Rèn luyện tư duy thiết kế phần mềm định hướng an toàn, bảo vệ dữ liệu cá nhân và quản lý phiên làm việc hiệu quả.
\end{itemize}

\section{Đối tượng và phạm vi nghiên cứu}

\subsection{Đối tượng nghiên cứu}
\begin{itemize}
  \item Kiến trúc tác nhân thông minh (Agent) và vòng lặp tác nhân (agent loop) ReAct.
  \item Cơ chế tối ưu ngữ cảnh và các phương pháp nén thông tin hội thoại.
  \item Khung công cụ (\texttt{Tool Harness}) và cơ chế gọi hàm (tool call).
  \item Quy trình trích xuất quy trình làm việc có kiểm soát (Reusable Workflow Extraction - HITL).
  \item Bộ lập lịch sự kiện chạy ngầm và cơ chế gộp tin nhắn (message coalescing).
\end{itemize}

\subsection{Phạm vi nghiên cứu}
Đề tài tập trung vào các nội dung thực tiễn sau:
\begin{itemize}
  \item Sử dụng môi trường Node.js và TypeScript để phát triển hệ thống \cite{nodejsDocs,typescriptDocs}.
  \item Sử dụng LangGraph làm động cơ điều phối trạng thái, grammY xây dựng kênh Telegram Bot và nền tảng LangSmith để giám sát, gỡ lỗi \cite{langgraphDocs,grammyDocs,langsmithDocs}.
  \item Triển khai tác nhân ưu tiên cục bộ (local-first), hỗ trợ gọi LLM cục bộ (Ollama) và kết nối với các mô hình ngôn ngữ lớn qua API đám mây thương mại \cite{ollamaDocs}.
  \item Lưu trữ an toàn dữ liệu người dùng cục bộ qua các tệp định dạng JSON gồm \texttt{reminders.json}, \texttt{memory.json} và \texttt{profile.json}.
\end{itemize}

Đề tài không thực hiện:
\begin{itemize}
  \item Tái huấn luyện hoặc tinh chỉnh (fine-tuning) các mô hình ngôn ngữ lớn từ đầu.
  \item Xây dựng giao diện Web Dashboard tự phát triển phức tạp hoặc ứng dụng di động native.
  \item Tích hợp các cổng thanh toán hoặc nền tảng đám mây khác ngoài Telegram và LangSmith.
\end{itemize}

\section{Phương pháp nghiên cứu}

\subsection{Nhóm phương pháp nghiên cứu lý thuyết và tài liệu}
\begin{itemize}
  \item Tìm hiểu các mô hình bộ nhớ tác nhân, đặc biệt là cơ chế quản lý bộ nhớ hợp nhất AgeMem \cite{yu2026agentic} hay hệ thống lưu trữ MemMachine \cite{memverge2026memmachine}.
  \item Nghiên cứu tiến trình phát triển cơ chế gọi công cụ của LLM Agent, từ gọi công cụ đơn lẻ đến điều phối đa công cụ phức tạp \cite{xu2026evolution}.
  \item Phân tích lý thuyết về thiết kế khung công cụ cho tác nhân qua mô hình thiết kế khung công cụ tác nhân (Agentic Harness Engineering - AHE) \cite{ahe2026agentic} và kiến trúc khung công cụ mở rộng (Extensible Harness) \cite{lastharness2026last}.
  \item Nghiên cứu ứng dụng mô hình tự thích ứng MAPE-K để giám sát và tối ưu hóa vòng lặp tác nhân \cite{emergentmindmape}.
  \item Khảo sát các mã nguồn mở thiết kế tác nhân cục bộ tối giản (như Nanobot \cite{hkudsnanobot} và OpenClaw \cite{openclawproject}) làm cơ sở định hình kiến trúc hệ thống.
\end{itemize}

\subsection{Nhóm phương pháp nghiên cứu thực nghiệm}
\begin{itemize}
  \item Xây dựng nguyên mẫu hệ thống MiniClaw chạy thử nghiệm thực tế.
  \item Đo lường và đánh giá các chỉ số: độ chính xác khi gọi công cụ, lượng token tiết kiệm sau khi nén ngữ cảnh, và thời gian phản hồi của hệ thống.
  \item Thu thập phản hồi từ người dùng thực nghiệm để điều chỉnh và tối ưu hóa quy trình tương tác.
\end{itemize}

\section{Cấu trúc đồ án}
Báo cáo đồ án được cấu trúc thành 6 chương chính thể hiện toàn bộ nội dung nghiên cứu và triển khai thực nghiệm hệ thống trợ lý ảo:
\begin{itemize}
  \item \textbf{Chương I: Tổng quan}: Trình bày lý do chọn đề tài, mục tiêu, đối tượng, phạm vi và phương pháp nghiên cứu.
  \item \textbf{Chương II: Cơ sở lý thuyết}: Giới thiệu các khái niệm về LLM, ReAct Agent, quản lý bộ nhớ tác nhân, tối ưu ngữ cảnh, khung công cụ và các công nghệ sử dụng.
  \item \textbf{Chương III: Tổng quan về hệ thống}: Trình bày ý tưởng thiết kế, các nguyên tắc cốt lõi như gộp tin nhắn (message coalescing), bảo mật ranh giới và vòng lặp hoạt động chính.
  \item \textbf{Chương IV: Phân tích và thiết kế hệ thống}: Chi tiết sơ đồ ca sử dụng, sơ đồ hoạt động, sơ đồ trình tự, thiết kế ERD, sơ đồ lớp, kiến trúc phân lớp và kịch bản tương tác người dùng.
  \item \textbf{Chương V: Cài đặt và kiểm thử}: Mô tả môi trường phát triển, cài đặt các thành phần chính và phân tích kết quả kiểm thử qua 5 nhóm kịch bản.
  \item \textbf{Chương VI: Kết luận}: Tổng kết kết quả đạt được, các hạn chế còn tồn tại và định hướng phát triển trong tương lai.
\end{itemize}
